\cotefile{91}
\scanpage{19}
marqués de $(\unsure{+|\star|\ast})$, dont les images dans le coin droit du bas sont les mêmes, on cherche un élément du coin gauche du haut qui donne ces éléments marqués. On note qu'il suffit pour ceci que $P(u_0,S)\ne\varnothing$, car alors le diagramme $(1)$ est isomorphe à $(2)$, or $(1)$ est \uline{cartésien}, grâce au fait que $M_{v_0}\to M_{u_0}$ est surjectif. Or l'obstruction à trouver un élément dans $P(u_0,S)$ est \deleted{\unsure{dans}}

\[
\partial(u_0)=\xi\in\operatorname{Ext}^{1}_{\unsure{\mathcal O_{P_0}|\underline{\mathcal O}_{P_0}}}
(P_0;I_{\Gamma_0},\mathcal O_{\Gamma_0}\otimes_{S_0}J)
\]

% slanted

\margin{\deleted{\unsure{sous réserve|sauf erreur} qu'il ne faut pas remplacer $P_0$ par $P_1$, \unsure{évident que|évident par}}}

où $P_0=X_0\times_{S_0}Y_0$, $\deleted{\unsure{\xi|\Gamma}}\Gamma_0$ est le sous-préschéma de $P_0$ graphe de $u_0$, et où $I_{\Gamma_0}$ est l'Idéal de $P_0$ qui le définit. C'est un module sur $A_0$, qui s'annule quand on \deleted{\unsure{tensorise par $A_{0\eta}$}}\inserted{localise en $\eta$}, donc Si $\pi$ est un paramètre de $A$, il existe un $n\geq 0$ tel que $\pi^n\xi=0$. Utilisant la suite exacte

\[
0\to J\xrightarrow{\pi^{\unsure{n|n+1}}}J\to J_n\to 0 \tag*{$(\times)$}
\]

et la suite correspondante

\[
0\to\mathcal O_{\Gamma_0}\otimes_{S_0}J\xrightarrow{\pi^n}\underline{\mathcal O}_{\Gamma_0}\otimes_{S_0}J\to\mathcal O_{\unsure{\Gamma_0|\Gamma_s}}\otimes_{S_0}J_n\to0
\]

et la suite exacte des $\operatorname{Ext}^{i}_{\mathcal O_{P_0}}(P_0;I_{\Gamma_0},-)$, on trouve que $\xi$ provient d'un élément

\[
\begin{aligned}
&\xi'\in\operatorname{Hom}_{\mathcal O_{P_0}}(P_0;I_{\Gamma_0},\underline{\mathcal O}_{\Gamma_0}\otimes_{S_0}J_n)\\
&\hspace{4em}\bmod\operatorname{Im}\operatorname{Hom}_{\mathcal O_{P_0}}(P_0;I_{\Gamma_0},\mathcal O_{\Gamma_0}\otimes_{S_0}J]
\end{aligned}
\]

% archive sheet number: 18
\scanpage{20}
Or le \inserted{premier terme du} deuxième \unsure{mb|m.b} s'écrit

\[
\begin{aligned}
\operatorname{Hom}_{\unsure{\mathcal O_{\Gamma_0}|\underline{\mathcal O}_{\Gamma_0}|\mathcal O_{P_0}}}\deleted{\unsure{\mathcal O_{P_0}\otimes}}(I_{\Gamma_0}/I_{\Gamma_0}^2,\mathcal O_{\Gamma_0}\otimes_{S_0}J_n)&\simeq\operatorname{Hom}_{\mathcal O_{X_0}}(u_0^*(\Omega^1_{Y_0/S_0}),\mathcal O_{X_0}\otimes_{S_0}\unsure{J_n|J}\unsure{)} \\
\deleted{\unsure{\text{et}}}\qquad &\simeq\operatorname{Hom}_{S_0}(M_{u_0},J_n)
\end{aligned}
\]

et de même le deuxième terme s'écrit $\operatorname{Hom}_{S_0}(M_{u_0},J)$, de sorte que l'obstruction peut être \unsure{considéré|considérée} comme un élément

\[
\xi''\in\operatorname{Hom}_{A_0}(M_{u_0},J_n)/\operatorname{Im}\operatorname{Hom}_{A_0}(M_{u_0},J\unsure{.})\hookrightarrow\operatorname{Ext}^{1}_{A_0}(M_u,J)
\]

[en utilisant encore, pour la dernière inclusion, la suite exacte $\times$]. Donc on obtient un élément \unsure{.}

\[
\xi''\in\operatorname{Ext}^{1}_{A_0}(M_u,J)
\]

annulé par une puissance de $\pi$, dont l'annulation est néc. et suffisante pour la possibilité de prolonger $u_0$ en $u:X\longrightarrow Y$. On trouve de même un

\[
\unsure{y|\eta|\gamma}''\in\operatorname{Ext}^{1}_{A_0}(M_{\unsure{v|v_0}},J)
\]

et nous admettrons que $y''$ est l'image de $\xi''$ grâce à $M_v\longrightarrow M_u$. Or \uline{$y''$ est nul}, car $v_0$ se prolonge en $v$, \deleted{\unsure{donc}} Utilisant la suite exacte

\[
0\to N\to M_{v_0}\to M_{u_0}\to0
\]

il vient, par la suite exacte des $\operatorname{Ext}^{i}(\,,J)$, que $\xi''$ correspond à un élément

% archive sheet number: 19
\scanpage{21}
\[
\xi'''\in\operatorname{Hom}(N,J)/\operatorname{Im}\operatorname{Hom}(M_{v_0},J)
\]

\deleted{Cet} Cet élément localisé en $\eta$ est nul, \deleted{\unsure{car}} et c'est là tout ce que nous pouvons tirer du fait que $u_0\circ\varphi_0$ et $u_{0\eta}$ se prolongent. Mais \uline{nous savons plus} : les prolongements \inserted{$v$ et $u_\eta$} peuvent se faire de façon que $v_\eta=\deleted{\unsure{\varphi_\eta}}u_\eta\varphi_\eta$. À priori, \deleted{\unsure{l'obs|l'idée}} une fois acquis le fait de la possibilité \uline{séparée} des deux obstructions, i.e. dans $(1)$ de l'existence de deux \unsure{élémts|éléments} dans les coins, la \deleted{\unsure{obs}} possibilité de les choisir de façon cohérente se trouve dans

\[
\begin{aligned}
&\operatorname{Hom}(M_{v_0},J_\eta)/\operatorname{Im}\big(\operatorname{Hom}(M_{v_0},J)+\operatorname{Hom}(M_{u_0},J_\eta)\big)\\
&\subset\operatorname{Hom}(N,J_\eta)/\operatorname{Im}\operatorname{Hom}(M_{v_0},J)
\end{aligned}
\]

Soit $\unsure{\gamma|y|\eta}$ cet élément. Notons que $\xi'''$ définit également un élément du dit quotient, puisque

\[
\operatorname{Hom}(N,J_{\deleted{\unsure{\eta|n}}})/\operatorname{Im}\operatorname{Hom}(M_{v_0},J)\subset\operatorname{Hom}(N,J_\eta)/\operatorname{Im}\ldots
\]

Nous admettrons que l'on a $\gamma=\xi'''$. Comme par hypothèse $\unsure{\gamma|\eta}'=0$, on a $\xi'''=0$, OK.

% archive sheet number: 20
\scanpage{22}
\uline{Pb} Simplifier \unsure{ce|le|les} \inserted{dernier} raisonnement, en éliminant les hypothèses de propreté sur $X,Z$ et le résultat délicat de l'existence de $M_u,M_v$, en \unsure{écrivant|énonçant|imposant} l'hypothèse de rigidité par

\margin{?}

\[
\operatorname{Hom}(u_0^*(\Omega^1_{Y_0/S_0}),\unsure{\mathcal O_{X_0}|\underline{\mathcal O}_{X_0}}\otimes G)\longrightarrow\operatorname{Hom}(v_0^*(\Omega^1_{X_0/S_0}),\mathcal O_{Z_0}\otimes G)
\]

injectif pour tt $G$ quasi-cohérent sur $S_0$.

\noindent\rule{\linewidth}{0.3pt}

% Two large sets of diagonal cancelling strokes cover the following paragraphs and display.

\deleted{Pour le \deleted{\unsure{th|Pb}}\inserted{critère} \uline{générique}, on est ramené à ceci. $S$ étant irréductible de pt générique $\eta$, $\unsure{S_\eta|S}$ muni d'un Idéal $J$ de carré nul, (ou $=$, si on y tient, annulé par le \unsure{nilradical}), d'où $S_0=V(J)$, et $u_0:X_0\to Y_0$ avec $v_0=u_0\varphi_0$ étant tels qu'on \uline{peut} \uline{prolonger \inserted{\unsure{de tout|d'abord}} en pt générique avec cette même condition}, prouver que c'est \unsure{possible|permis} en \deleted{\unsure{les autres}}\inserted{tous les} pts d'un voisinage.}

\deleted{Or, on a dit que \deleted{l'obstruction à relever était}\inserted{\unsure{dès|ds} la question de relever} à un $S'_{0s}\supset S_{0s}$ donné, \deleted{\unsure{était indépendante de la question}}\inserted{la présence de la} condition imposée par $v$ était inessentielle, donc la question est \deleted{\unsure{de donner à la}}\deleted{\unsure{la question}}\deleted{\unsure{par la nullité de}} d'une certaine obstruction}

\deleted{\[
\xi\in\operatorname{Ext}^{1}_{\mathcal O_{P_0}}(P_0;I_{\Gamma_0},\mathcal O_{\Gamma_0}\otimes_{S_0}J)
\]}

% archive sheet number: 21
\scanpage{23}
$S$ \inserted{\unsure{nœth|noeth}} muni d'un Idéal \inserted{\deleted{\unsure{qu}}coh.} de carré nul $J$,

$Z\xrightarrow{\unsure{\varphi,v|v}}X;Y$ sur $S$ d'où $v_0:Z_0\mathbin{-}X_0,Y_0$ sur $S_0$

\[
S_0=V(J)
\]

$u_0:X_0\to Y_0$ tel que $u_0\varphi_0=v_0$

% slanted insertion in left margin, with a leader to the equation below
On va étudier l'obstruction à relever $u_0$ en $u:X\to Y$ tel que $u\varphi=v$\inserted{ et l'obstruction analogue pour la situation déduite à l'aide d'un hom \unsure{$u$|$w$}$:J\to G$, ou après localisation\ldots} À priori, \deleted{l'obstruction à relever}\inserted{les relèvements de} $u_0$ en un $u$ (sans condition supplémentaire imposée) correspondent aux splittings d'une extension sur $P_0$

\[
0\mathbin{-}\unsure{\mathcal O_{\Gamma_0}|\underline{\mathcal O}_{\Gamma_0}}\otimes_{S_0}J\longrightarrow E\longrightarrow\underline I_{\Gamma_0}\mathbin{-}0
\]

% slanted

\margin{SGA 3 \uline{III} 4.2}

les relèvts pour $v_0$ correspondent aux splittings d'une extension \deleted{\unsure{de}} sur $P'_0=Z_0\times_{\unsure{S_0|S}}Y_0$

\[
0\mathbin{-}\mathcal O_{\Gamma'_0}\otimes_{S_0}J\mathbin{-}E'\mathbin{-}\underline I_{\Gamma'_0}\mathbin{-}0
\]

D'ailleurs, on constate qu'on a un homomorphisme d'extensions, compatible avec $P'_0\to P_0$

\diagram{\begin{tikzcd}[ampersand replacement=\&,column sep=small]
0\arrow[r] \& \mathcal O_{\Gamma_0}\otimes_{S_0}J\arrow[r]\arrow[d] \& E\arrow[r]\arrow[d] \& I_{\Gamma_0}\arrow[r]\arrow[d] \& 0\\
0\arrow[r,dash] \& \mathcal O_{\Gamma'_0}\otimes_{S_0}J\arrow[r,dash] \& E'\arrow[r] \& \underline I_{\Gamma'_0}\arrow[r] \& 0
\end{tikzcd}}

(en fait, $E'$ se déduit de $E$ par extension de la base $P'_{\unsure{\xi|0}}\mathbin{-}P_{\unsure{\xi|0}}$ \unsure{ou} avec

\diagram{\begin{tikzcd}[ampersand replacement=\&,column sep=small,row sep=small]
\& P_0\arrow[d]\\ Z_0\arrow[r] \& X_0
\end{tikzcd}}

\unsure{;]|;} il faut préciser ce pt \ldots{} ]. \deleted{\unsure{Considérons, pour un|Conditions, puis un}}\deleted{\unsure{$G$ quasi-cohérent sur}}\inserted{D'ailleurs, la} deuxième ligne est munie d'un splitting, et les solutions cherchées sont les splittings de la première \uline{compatibles} avec celui donné pour la seconde.

% archive sheet number: 22
\scanpage{24}
Ceci posé, regardons, pour un $G$ quasi-cohérent variable sur $S_0$, les $q_*\underline{\operatorname{Hom}}(\,,\unsure{\mathcal O_{\Gamma_0}|\underline{\mathcal O}_{\Gamma_0}}\otimes_{S_0}G)$, $q'_*\underline{\operatorname{Hom}}(\,,\mathcal O_{\Gamma'_0}\otimes_{S_0}G)$ relatives aux deux lignes, on trouve \inserted{(via les Modules représentatifs)} un diagramme commutatif de \deleted{suites exactes}

\diagram{\begin{tikzcd}[ampersand replacement=\&,column sep=normal]
\& \& L\arrow[r,"i"] \& M\arrow[r,"d"] \& N\arrow[r] \& 0\\
(D) \& 0\arrow[r] \& L'\arrow[u,"\alpha"']\arrow[r,"i'"] \& M'\arrow[u,"\beta"']\arrow[r,"d'"] \& N'\arrow[u,"\gamma"']\arrow[r,dash] \& 0\,, 
\end{tikzcd}}

de plus on a un hom can. $\sigma$

\diagram{\begin{tikzcd}[ampersand replacement=\&,row sep=small,column sep=small,execute at end picture={\draw ([xshift=-4pt,yshift=5pt]js.north west) -- ([xshift=-8pt,yshift=4pt]ls.north west) -- ([xshift=5pt,yshift=1pt]ls.north east); \draw ([xshift=-4pt,yshift=0pt]js.south west) -- ([yshift=-6pt]js.south) -- ([xshift=5pt,yshift=-1pt]ls.south east);}]
\& \& |[alias=ls]|L\arrow[dl,"\sigma"']\\
\& |[alias=js]|J\arrow[dl,dashed,"\lambda"'] \& L'\arrow[u,"\alpha"']\\
G
\end{tikzcd}}

\diagram{\begin{tikzcd}[ampersand replacement=\&]
L'\&M'\arrow[l,"\pi'"']
\end{tikzcd}}

correspondant à l'identité $\mathcal O_{\Gamma_0}\otimes_{S_0}J\longrightarrow\mathcal O_{\Gamma_0}\otimes_{S_0}\underset{\substack{\Vert\\G}}{J}$, et un \uline{splitting} \inserted{$M'\xrightarrow{\unsure{\pi|\pi'}}L'$} de la deuxième ligne.

Ceci dit, pour un $\lambda:J\to G$ donné, les solutions du pb relativement à $G$ sont les hom $M\xrightarrow{\mu}G$ qui relèvent \inserted{i.e. $\boxed{\mu i=\lambda\sigma}$} le composé $L\xrightarrow{\lambda\sigma}G$, et tels que $\mu\beta=(\underbrace{\lambda\sigma}\alpha)\pi'$. \deleted{\unsure{Il}} Pour qu'il en existe un, il f. et s. que

\margin{Condition de non ramification en $u_0$}

$1^\circ)$ $(\lambda\sigma)(\operatorname{Ker}i)=0$ i.e. $\lambda(R_1)=0$, où $R_1=\sigma(\operatorname{Ker}i)\subset J$

\margin{Alors on peut remplacer $L$ par $L/\operatorname{Ker}i$}

[\deleted{i.e. $\lambda$ se factorise en $\bar\lambda:\bar J=J/K\longrightarrow G$}].

$2^\circ)$ [Tenant compte que \uline{$\gamma$ \deleted{\unsure{est}}surjectif} ] $\lambda\sigma\alpha\pi'$ \unsure{vaut|nulle} sur \deleted{\unsure{$R'$}}$\operatorname{Ker}\beta$, i.e. $\lambda(R_2)=0$, où $R_2=\unsure{\lambda}\sigma\alpha\pi'(\operatorname{Ker}\beta)$

% archive sheet number: 23
\scanpage{25}
\diagram{\begin{tikzcd}[ampersand replacement=\&,column sep=normal,row sep=small]
\operatorname{Ker}i\arrow[d,very thick]\arrow[r,dashed] \& L\arrow[dl,dashed] \&\\
J \& L'\arrow[l,dashed]\arrow[u,dashed,"\unsure{2|\alpha|?}"'] \& M'\arrow[l,dashed]\\
\& \& \operatorname{Ker}\beta\arrow[u,dashed]\arrow[ull,thick]
\end{tikzcd}}

\diagram{\begin{tikzcd}[ampersand replacement=\&,row sep=small]
L+M'\arrow[r,very thick,"{(i,\beta)}","{\text{surj}}"' ]\arrow[d,"{(\sigma,\sigma\alpha\pi')}"',"{=\rho}"] \& M\arrow[ddl,dashed,bend left=12,"\mu"]\\
J\arrow[d,"\lambda"'] \&\\
G \&
\end{tikzcd}}

\margin{condition de non ramification}

\margin{le foncteur en $G$ est repr. par $\operatorname{Coker}\rho$, dont la formation commute à tout \unsure{chgt|chgt.} de base}

En résumé, la condition en question est l'annulation de $\lambda$ sur $R=R_1+R_2$, donc la factorisation de $\lambda$ par $J/R\longrightarrow G$.

D'ailleurs, la formation du quotient ainsi \unsure{obtenu|obtenus} $J/R$ de $J$ commute à l'extension plate de la base $S$ [en particulier à la localisation].

Ceci assure la validité des conditions \uline{7$^\circ$} et \uline{8$^\circ$}.

\noindent\rule{\linewidth}{0.3pt}

\uline{Corollaires du \unsure{th|Th}}

\uline{Corollaire 1} Soient $P/X/k$ sur $k$ alg. clos \inserted{(\uline{sép. clos suffit}!)} $X$ propre sur $k$, $P/X$ loc de \deleted{\unsure{présentation}}\inserted{type} fini. Alors $\prod_{X/k}P/\deleted{X}$ représentable par préschéma loc \deleted{\unsure{de}} algébrique.

Si $P/X$ quasi-projectif\inserted{ ou seulmt à fibres \unsure{qs. pro.|qu.-pr.}}, alors $\prod_{X/k}\longrightarrow{}P/k$ est essentiellement quasi-projectif [i.e. tt ouvert quasi-compact est quasi-projectif] sur $k$.

\unsure{$\swarrow$} il suffit que pour tt $x\in\operatorname{Ass}\mathcal O_X$, $\exists\,y\in\bar x$ tel que la fibre $P_y$ soit quasi-projective.

% archive sheet number: 24
\scanpage{26}
\unsure{Corollaire|Corollaires} 1 à 7 contenu\unsure{s|{}} dans le suivant :

\uline{Théorème}

$X$ propre et plat de prés. finie sur $S$ \deleted{\unsure{localement}}, $P/X$ loc. de présentation finie sur $X$ $Z_1\ldots Z_n$ des sous-préschémas \inserted{fermés} de $X$, plats et de présentation finie sur $S$. On suppose

(a) Pour tt $i$, $\prod_{Z_i/S}P_{Z_i}/Z_i$ représentables [p. ex. $Z_i/S$ radiciels, ou $Z_i/S$ et $P_{Z_i}/Z_i$ projectifs]

(b) Pour tout $i$, $X$ est normalement plat le long de $Z_i$ [p. ex. $Z_i$ simple $/S$ et $X/S$ simple aux pts de $Z_i$]

(c) Pour tout $s\in S$, et tt $x\in\operatorname{Ass}\unsure{\mathcal O_{X_{\bar s}}|\underline{\mathcal O}_{X_{\bar s}}}$, il existe $i$ tel que $Z_{i\bar s}\cap\bar x\ne\varnothing$, i.e. $(\bigcup Z_i)_{\bar s}\cap\bar x\ne\varnothing$ $[\bar s=\operatorname{Spec}\overline{k(s)}]$

Sous ces conditions, (i) $\prod_{X/S}P/X$ est représentable par un préschéma loc de \deleted{type}\inserted{présentation} finie et séparé sur $S$

\deleted{\unsure{Toute partie}} (ii) \inserted{Suppos $S$ est quasi-séparé} Pour \unsure{tte|tt.} partie \inserted{ouverte} quasi-compacte $U$ de $\prod_{X/S}P/X$, \deleted{\unsure{il existe un entier $n$}}\deleted{\unsure{tel que}} le morphisme $U\longrightarrow\prod_{1\leq i\leq n}\deleted{\unsure{\underline{\operatorname{Hom}}(Z}}\prod_{Z_i/S}P_{Z_i}/Z_i$ est quasi-affine

(iii) Par suite, si les $\prod_{Z_i/S}P_{Z_i}/Z_i$ sont quasi-projectifs \inserted{sur $S$ \deleted{\unsure{quasi-sép}}}, p. ex. $Z_i/S$ fini et $P_{Z_i}/Z_i$ quasi-projectif, alors toute partie quasi-compacte de $\prod_{X/S}P/X$ est quasi-projective sur $S$ [de façon précise, il existe un Module inversible $\underline L$ sur $\prod_{X/S}P/X$ qui est essentiellement ample rel $/S$].

% archive sheet number: 25
\scanpage{27}
\uline{Pb} Les composantes connexes de $\prod_{\deleted{\unsure{X}}X/k}\unsure{P'|P}/X$ sont elles de type fini sur $k$ ?

\uline{Corollaire 2} Supposons, \deleted{\unsure{que}}\inserted{pour} $P/X/k$, $X$ propre sur $k$, $P/X$ loc. de type fini \deleted{\unsure{et diff. simple}}, $k$ corps (pas néc. sép. clos) que \deleted{tt} pour tt $x_i\in\operatorname{Ass}\unsure{\mathcal O_X|\underline{\mathcal O}_X}$, $\exists\,x'_i\in\bar x$ tel que $k(x'_i)/k$ radiciel. \inserted{[plus gén, tel que $\prod_{k(x_i)/k}P_{x_i}/k(x_i)$ représentable]} Alors $\prod_{X/k}P/X$ représentable.

% slanted

\margin{\deleted{\unsure{diff$^{t}$ simple de la} \unsure{projection aux germes}, car le foncteur ne change pas alors en se bornant à l'ouvert $P'$ de $P$ où $P$ est \uline{simple} sur $X$}}

\margin{Si les $P_{x'_i}$ sont quasi-projectifs, alors $\prod_{X/k}P/X$ est essentiellement quasi-projectif}

P. ex. si $X$ \deleted{est} est \deleted{\unsure{irréductible}}\deleted{\unsure{irréductible}}, \inserted{\unsure{irrédondant|irrédundant}} et admet un pt $x$ radiciel sur $k$ \deleted{\unsure{et}}, $\prod_{X/k}P/X$ existe.

% Tall right bracket along the three lines of this corollary.
\margin{??}
\uline{Corollaire 3} $P/X/k$ comme dessus. Alors $\exists$ $k'/k$ extension \uuline{finie} \uline{séparable}, telle que $\prod_{X'/k'}P'/X'$ représentable.

% slanted

\margin{\deleted{Détailler la démonstration, lorsqu'on ne suppose pas $P/X$ diff. simple, mais $P/X$ de type fini : l'\uline{énoncé} est-il \uline{vrai} ???}}

\uline{Corollaire 4} \deleted{\unsure{Si}} $P/X$ à fibres quasi-projectives, et $X/k$ propre \deleted{\unsure{plat}}, alors $\prod_{X/k}P/X$ \unsure{exist|existe} et est ess. quasi-proj. (\unsure{Varie|Voir} : Cor. 1)

\uline{Corollaire 5} $P/X/S$, $X$ propre et plat sur $S$, \deleted{\unsure{un}}\unsure{nœth|noeth} intègre, $P/X$ loc. de type fini. Alors $\exists$ \deleted{voisinage} ouvert $U\ne\varnothing$, et morphisme fini étale surjectif $U'\to U$, tel que $\prod_{X_{U'}/U'}P_{U'}/X_{U'}$ représentable.

% slanted

\margin{\deleted{\deleted{\illegible{1}} Faut-il supposer $P/X$ diff. simple ??}}

Si $X_\eta/k(\eta)$ satisfait la condition énoncée dans Cor. 2, on peut même prendre $U'=U$

% archive sheet number: 26
\scanpage{28}
\uline{Corollaire \unsure{5|6}} $P/X/S$, avec $X\longrightarrow S$ propre plat à fibres géométriques intègres. Supposons qu'il existe une section $\sigma$ de $X/S$ le long de laquelle $X/S$ soit \deleted{simple}\inserted{lisse}\deleted{ et $P/X$ \unsure{est partout} diff. simple}. Alors $\prod_{X/S}P/X$ représentable. Si $P\times_X(S,\sigma)$ est quasi-projectif sur $S$, alors $\prod_{X/S}P/X$ est essentiellement quasi-projectif

\uline{Corollaire 7} $P/X/S$, avec $X\to S$ simple propre \deleted{\unsure{connexe|Comme dessus}, $P/X$ diff simple} \deleted{à fibres géométriques intègres} Alors $\prod_{P/X}X/S$ est localement représentable pour la top. étale.

Remplaçant la top étale par la top. fpqf (fid plat quasi finie), le résultat \unsure{6|/} \unsure{devrait|doit|demeure} être valable sans supposer $X\to S$ simple, mais seulemt propre, plat, à fibres \deleted{géométriques}\inserted{sans cycles} associés immergés, par essentiellement la $=$ démonstration.

\uline{Corollaire 8} $X$ et $Y$ propres sur \deleted{\unsure{$k$}} corps $k$. Alors $\underline{\operatorname{Isom}}_k(X,Y)$ est représentable. En particulier, $\underline{\operatorname{Aut}}_k(X)$ est représentable

% archive sheet number: 27
\scanpage{29}
\uline{Corollaire 9} Soit $X/S$ propre et plat sur $S$ \unsure{noethérien|nœthérien} intègre. Alors $\exists$ $U$ non vide, \deleted{et $U'\longrightarrow U$ fini, \unsure{fidèle|fid. plat}, surjectif}\inserted{et $U'\to U$ fini étale \deleted{\unsure{fidèle}} surjectif} tel que $\underline{\operatorname{Aut}}_S(X)\times_S U'$ soit représentable.

Il est plausible qu'on puisse $=$ prendre $U'=U$, c'est lié à une question d'effectivité des données de descente\ldots.

% slanted

\margin{cf Raynaud}

% archive sheet number: 28
\scanpage{31}
\handwritten{11.9.1967}\\

Dear Anantharaman,\\

Matsumura proved that if $X$ is proper over a field $k$, then $\underline{\operatorname{Aut}}_{X/k}$ is representable by a group scheme locally of finite type over $k$. I think I can systematize the key step of his argument in the following way. Consider a scheme $S$, and a morphism

\[
\phi: Z\handwritten{\longrightarrow} X
\]

of $S$-schémes which are proper, flat and of finite presentation. Let $Y$ be locally of finite présentation and separated over $S$, then $\phi$ induces a homomorphisme of functors $u\handwritten{\leadsto}u\phi$ ;

\[
\phi':\underline{\operatorname{Hom}}_S(X,Y)\handwritten{\longrightarrow}\underline{\operatorname{Hom}}_S(Z,Y)\quad .
\]

Then one can défine a subfunctor of $\underline{\operatorname{Hom}}_S(\deleted{\unsure{Z}}X,Y)$ where $\phi'$ is ``unramified'' in a rather obvious sense, and this turns out to be an ``open subfunctor'', say $\underline{\operatorname{Hom}}_S(X,Y;\phi)$ . Now look at the induced homomorphism

\[
\underline{\operatorname{Hom}}_S(X,Y;\phi)\handwritten{\longrightarrow}\underline{\operatorname{Hom}}_S(Z,Y)\,.
\]

Using the main result of Murre's talk, one can prove that the latter morphism is representable by unramified separated morphisms locally of finite presentation; as a consequence, if $\underline{\operatorname{Hom}}_S(Z,Y)$ is representable, so is $\underline{\operatorname{Hom}}_S(X,Y;\phi)$.

To get, given $X$ and $Y$, a representability theorem for $\underline{\operatorname{Hom}}_S(X,Y)$, on\deleted{\unsure{ye|es|ge}}\unsure{e} tries to find morphisms $\phi_i:Z_i\handwritten{\longrightarrow}X$ as above, such that the open subfunc\deleted{\unsure{ot}}tors $\underline{\operatorname{Hom}}_S(X,Y;\phi_i)$ cover $\underline{\operatorname{Hom}}_S(X,Y)$ (as a fpqc sheaf), and such that the functor\deleted{\unsure{z}}s $\underline{\operatorname{Hom}}_S(\deleted{\unsure{X}}Z_i,Y)$ are all representable. If for instance $\deleted{\unsure{X}}S$ is the spectrum of a field $k$, and if $X$ has ``enough'' points radicial over $k$ (which is always true if $k$ is alg. closed) then we can take for $Z_i$ all finite subs\deleted{\unsure{h}}\unsure{c}hemes of $X$ whose points are radicial over $k$, and we get

% archive sheet number: 30
\scanpage{32}
that $\underline{\operatorname{Hom}}_S(X,Y)$ is representable (ahy $Y$ locally of finite presentation \inserted{\handwritten{and separated}} over $k$); if we do not make any assumption on\unsure{'|{}} $X$ except properness over $k$, the previous assumption becomes true after finite ground-field extension $k'/k$, so that we get that for every $Y$ as above, $\underline{\operatorname{Hom}}_{\deleted{\unsure{s}}S}(X,Y)\mathbin{\mathrm{x}}_S\operatorname{Spec}(k')$ is representable. From this \unsure{Matsumuràs|Matsumura's|Matsumuras} theorem stated at the beginning follows in a standart way by descebt arguments.\unsure{'|{}} T\deleted{\unsure{n}}he result holds too for $\underline{\operatorname{Isom}}_k(X,Y)$ instead of $\underline{\operatorname{Aut}}_k(X)$, but as you probably know, $\underline{\operatorname{Hom}}_S(X,Y)$ is not always representable, even if $X$ is a quadratic extension of $S=\operatorname{Spec}k$, $Y$ being proper non projective.

Over an arbitrary base $S$, one can give a fairly general statement of a representability theorem,--the points \deleted{\unsure{rad}}radicial \unsure{over|ober} $k$ used above being replaced by suitable flat subschemes of $X$. As \deleted{applications}\inserted{particular cases}, we get for instance that if $X$ has integral \unsure{geomtric|geometric} fibers and a section\deleted{x} \inserted{alo\deleted{\unsure{g}}ng which $X$ is smooth,} then $\underline{\operatorname{Hom}}_S(X,Y)$ is representable; and if $X$ has reduced geometric fibers, then $\underline{\operatorname{Hom}}_S(X,Y)$ is repres\deleted{\unsure{a}}entable locally for the étale topology over $S$. Also, if $Y$ is quasi-projective over $S=\operatorname{Spec}k$, then $\underline{\operatorname{Hom}}_S(X,Y)$ is representable.

To gix the ideas, I gave the statements for $\underline{\operatorname{Hom}}_S(X,Y)$, but one has quite analoguous resultst of course for the $\handwritten{\prod}_{X/S}P/X$ fonctors, w\deleted{\unsure{h}}hich I guess will imply ra\unsure{hher|ther} formally the other ones\unsure{x}.

If you are interested, I can send you a photocopy of the statement of the general theorem of representability I alluded to\deleted{\unsure{,}} above, and a couple of coroliaries (I \unsure{already|alréady} listed here the most striking ones).

Sincerely your's\\

% archive sheet number: 31
